\section{Conclusion}
\label{sec:conclusion}

This study asked whether a support office should use a black-box model, the tabular foundation
model \tabicl, to prioritise capacity-limited student outreach, and on this dataset the answer is
negative. \tabicl ranked students well, but its differences from XGBoost, logistic regression and an
interpretable EBM lay within paired-bootstrap uncertainty, so its additional complexity is
unsupported. Its SHAP explanations were faithful under deletion, yet sensitive to the random seed
and comparison group, and they mostly restated weak first-semester results. In the human study,
charts did not improve predictions of the flag (difference $-0.016$) but slightly raised confidence,
and one unflagged case was misread more often with charts.

The EBM is consequently the reference model that a support office should validate next: it matched
the foundation model at the operating point, and its explanations are exact contributions of the
model itself rather than estimates that change on recomputation. A next-stage design should combine
a zero-pass check with this transparent ranking, display the cohort cut-off and the uncertainty of
each explanation, and keep an adviser responsible for every contact. As the evidence comes from one
institution and proxy users, evaluation on a current cohort with support staff is required before
any use, and the system must never be used to deny support.
